\documentclass[reqno]{amsart}
\usepackage{a4wide}

\AtBeginDocument{{\noindent\small
\emph{Electronic Journal of Differential Equations},
Vol. 2026 (2026), No. 19, pp. 1--16.\newline
ISSN: 1072-6691. URL: https://ejde.math.txstate.edu, DOI: 10.58997/ejde.2026.19}
\thanks{\copyright 2026. This work is licensed under a CC BY 4.0 license.}
\vspace{8mm}}

\begin{document}
\title[\hfilneg EJDE-2026/19\hfil Generalized quasilinear Schr\"odinger equations]
{Positive solutions for generalized quasilinear Schr\"odinger equations}

\author[X. Long, X.-D. Fang \hfil EJDE-2026/19\hfilneg]
{Xuan Long, Xiang-Dong Fang}

\address{Xuan Long\newline 
School of Mathematical Sciences, 
Dalian University of Technology, 
116024 Dalian, China}
\email{22401008@mail.dlut.edu.cn}

\address{Xiang-Dong Fang (corresponding author)\newline
School of Mathematical Sciences, 
Dalian University of Technology, 
116024 Dalian, China}
\email{fangxd0401@dlut.edu.cn}

\thanks{Submitted January 5, 2026. Published February 26, 2026.}
\subjclass[2020]{35J20, 35J62, 49J35}
\keywords{Quasilinear Schr\"odinger equation;  asymptotically linear; Nehari manifold}

\begin{abstract}
This article concerns the existence of solutions to the generalized quasilinear Schr\"odinger equation
\begin{equation*}
-\operatorname{div} (g^2(u)\nabla u)+g(u)g'(u)|\nabla u|^2+V(x)u=f(x,u), \quad 
x\in \mathbb{R}^{N},
\end{equation*} 
where  $N \ge 1$, $g \in C^1(\mathbb{R})$ and the nonlinearity is asymptotically linear at infinity. 
Through the stretching transformation, we adjust the approach in \cite{Ambrosetti} %[1]
and overcome the impediment that there might not exist a $t>0$ such that $t u_\infty(x-y)$ 
belongs to the Nahari manifold, where $u_\infty$ is a  ground state solution for the limiting problem.
\end{abstract}

\maketitle
\numberwithin{equation}{section}
\newtheorem{theorem}{Theorem}[section]
\newtheorem{lemma}[theorem]{Lemma}
\newtheorem{proposition}[theorem]{Proposition}
\newtheorem{remark}[theorem]{Remark}
\allowdisplaybreaks

\section{Introduction}

This article concerns the generalized quasilinear Schr\"odinger equation
\begin{equation}
\label{p1}
-\operatorname{div} (g^2(u)\nabla u)+g(u)g'(u)|\nabla u|^2+V(x)u=f(x,u), \quad x\in \mathbb{R}^{N},
\end{equation}
where $N \ge 1$, the limits of $V$ and $f$ exist as $|x| \to \infty$ and $g$  is continuously differentiable with a nonnegative derivative on $[0, \infty)$. 

Solutions to equation \eqref{p1} are directly associated with the stationary wave solutions of a 
time‑dependent quasilinear Schr\"odinger equation, which takes the form
\begin{equation}
\label{pp2}
i\Phi_t=-\Delta \Phi+W(x)\Phi-q(x,|\Phi|^2)\Phi-\Delta(l(|\Phi|^2))l'(|\Phi|^2)\Phi,
\end{equation}
where $\Phi(x,t)$ denotes the wave function, $W(x)$ stands for  the  potential, while $q$ and $l$  
are functions selected appropriately. 
This equation arises in the search for solitary wave solutions to a general quasilinear 
Schr\"odinger equation, which models several physical phenomena, and further details can be 
found in \cite{Wang2, Wang1, Poppenberg} for an explanation. 
By substituting the standing wave ansatz $\Phi(x,t)=e^{-i\lambda t}u(x)$ into equation \eqref{pp2}, 
one obtains the time‑independent elliptic equation
\begin{equation}
\label{p2}
-\Delta u+V(x)u-\Delta\bigl(k(u^2)\bigr)k'(u^2)u=f(x,u), \quad x\in \mathbb{R}^{N},
\end{equation}
which constitutes a particular case of problem \eqref{p1} with the coefficient function 
$g^2(u)=1+\frac{[(k(u^2))']^2}{2}$.

Two physically relevant specializations of this structure are noteworthy.
First, choosing $g^2(u):=1+2u^2$, which corresponds to $k(s)=s$, yields the superfluid film 
equation 
\begin{equation}
\label{p3}
-\Delta u+V(x)u-\Delta(u^{2})u=f(x,u), \quad x\in \mathbb{R}^{N}.
\end{equation}
Second, the choice $g^2(u):=1+\frac{u^2}{2(1+u^2)}$ leads to
\begin{equation}
\label{p4}
-\Delta u+V(x)u-\Bigl[\Delta\bigl(1+u^{2}\bigr)^{1/2}\Bigr]\frac{u}{2(1+u^{2})^{1/2}}=
f(x,u), \quad x\in \mathbb{R}^{N}.
\end{equation}
The latter model describes the propagation and self‑channeling of an intense ultrashort laser 
pulse inside a nonlinear optical medium.

There has been significant attention in the literature focused on  \eqref{p3}. 
The pioneering work in this direction was conducted by Poppenberg et al.\ \cite{Poppenberg}, leveraging variational methods. Subsequently, Wang \cite{Wang1} demonstrated that a positive 
solution exists, via a constrained minimization argument. A significant methodological advancement 
was made in \cite{Wang2}, which was achieved by reformulating the original quasilinear equation as 
a semilinear problem through an Orlicz-space-based variable transformation. Later, Colin et al.\ 
\cite{colin} introduced a modified transformation of variables, which allowed the analysis to be 
carried out within the conventional Sobolev space setting. Further extending these results, 
Silva et al.\ \cite{Silva} proved the presence of  nonnegative  solutions under asymptotically 
periodic and critical growth conditions.

Following these foundational works, considerable attention has been devoted to \eqref{p3}. 
For a broader overview of subsequent developments, we refer the reader to 
\cite{fang3, fang2, Jeanjean, Lehrer,  Xue} and the references cited therein. 

Our analysis focuses on the problem
\begin{equation}
\label{q1}
-\operatorname{div} (g^2(u)\nabla u)+g(u)g'(u)|\nabla u|^2+V(x)u=f(x,u), \quad u\in H^1(\mathbb{R}^{N}).
\end{equation}
We define   $E$ as the Sobolev space $H^{1}(\mathbb{R}^{N})$  with the norm 
$$
\|u\|:=\Big(\int_{\mathbb{R}^N}|\nabla u|^2+V(x)u^2\Big)^{1/2}.
$$ 
We set $G(t):=\int^t_0g(\tau)\,d\tau$,  $F(x,t):=\int^t_0f(x,\tau)\,d\tau$ and 
$F_{\infty}(t):=\int^t_0f_{\infty}(\tau)\,d\tau$. Our analysis is based on the following assumptions: concerning $V$, $g$, and $h$:
\begin{itemize}
\item[(A1)] $V$ is continuous and $0<\inf_{y\in{\mathbb{R}^N}} V(y)\leq V(x)\leq V_\infty$ for all
 $x\in {\mathbb{R}^N}$, where $V_{\infty}:=\lim_{|x| \to \infty} V(x)$.

\item[(A2)]  \(g\) is of class \(C^1\) on \(\mathbb{R}\), even, and  increasing on $ (0,\infty)$, with $g(0)=1$.

\item[(A3)] $f(x,t)\in C({\mathbb{R}^N}\times\mathbb{R}^+)$ and for every $t\geq 0$, 
$\lim_{|x|\to\infty}f(x,t)=f_{\infty}(t)$, where $f_{\infty}(t)$ is continuous in $\mathbb{R}^+$.

\item[(A4)] $\lim_{t\to0^+}\frac{f(x,t)}{t}=0$, uniformly for $x\in{\mathbb{R}^N}$.

\item[(A5)] $\lim_{t\to+\infty}\frac{f(x,t)}{g(t)G(t)}=s(x)$, uniformly for $x\in{\mathbb{R}^N}$, where $s(x)$ 
is continuous, $\lim_{|x| \to \infty} s(x) = s_{\infty}$ and $s(x)>\frac{V(x)}{g^2(\infty)}$ for every 
$x\in{\mathbb{R}^N}$, with $\frac{1}{g(\infty)}:=\lim_{t\to\infty}\frac{1}{g(t)}$ .

\item[(A6)] $t\mapsto \frac{f(x,t)}{g(t)G(t)}$ and $t\mapsto \frac{f_\infty(t)}{g(t)G(t)}$ are increasing 
on $(0,\infty)$. Moreover, on the set where $g(t)=1$, the aforementioned monotonicity is strict.

\item[(A7)] $\frac{f_{\infty}(t)-f(x,t)}{g(t)G(t)}\leq s_{\infty}-s(x)$, for all $(x,t)\in{\mathbb{R}^N}\times\mathbb{R}^+$.

\item[(A8)] If $g\equiv 1$, then $\lim_{t\to+\infty}\frac{1}{2}s(x)t^2-F(x,t)=\infty$, for a.e.\
 $x\in {\mathbb{R}^N}$.

\item[(A9)] For each $\delta>0$, $\lim_{|x|\to\infty}f(x,t)=f_{\infty}(t)$, uniformly for 
$0\leq t\leq\delta$.

\item[(A10)]   If $\omega_n \rightharpoonup 0$ in $E$, then for every $\gamma>0$, 
there exists $y\in{\mathbb{R}^N}$ such that 
\begin{gather*}
\liminf_{l\to \infty}\int_{|x-y|>l,\,|\omega_n|<\gamma}F(x,G^{-1}(|\omega_n|))-F_{\infty}(G^{-1}(|\omega_n|))=0, \\
\lim_{l\to \infty}\int_{|x-y|>l,\,|\omega_n|<\gamma}[\frac{f(x,G^{-1}(|\omega_n|))}{g(G^{-1}
(\omega_n))}-\frac{f_{\infty}(G^{-1}(|\omega_n|))}{g(G^{-1}(\omega_n))}]^2=0,
\end{gather*}
uniformly for $n$.

\item[(A11)] $0<s_{\infty}-\frac{V_{\infty}}{g^2(\infty)}\leq\inf _{x\in{\mathbb{R}^N}}[s(x)-\frac{V(x)}{g^2(\infty)}]
+\mu(\inf _{x\in{\mathbb{R}^N}}[s(x)-\frac{V(x)}{g^2(\infty)}])$, i.e., 
$\sup _{x\in{\mathbb{R}^N}}[\frac{V(x)-V_{\infty}}{g^2(\infty)} +s_{\infty}-s(x)]\leq\mu(\inf _{x\in{\mathbb{R}^N}}[s(x)-\frac{V(x)}{g^2(\infty)}])$, 
where $\mu(\cdot)$ is as in Lemma \ref{L3.17}.
\end{itemize}

\begin{theorem}\label{thm1.1}
Under assumptions {\rm (A1)--(A8), (A11)}   and at least one of {\rm (A9)} and {\rm (A10)},
we further impose the condition
\begin{itemize} %\label{H}
\item[(A12)] The minimal energy value $c_{\infty}$ corresponding to equation \eqref{q4} is an isolated critical level of the functional $J_{\infty}$.
\end{itemize}
Then \eqref{q1} admits a positive solution. 
\end{theorem}

\begin{remark}\label{r1} \rm
To prove that \eqref{q1} has a positive solution, we  can let  $f$ be odd for $t$.
\end{remark}

\begin{remark}\label{r2} \rm
For \( b > c_0=c_\infty \), we follow the approach in \cite{Ambrosetti}. 
Unlike \cite{Ambrosetti}, however, we cannot guarantee that \( u_\infty(x - y) \in \mathcal{T} \) 
for arbitrary \( y \) under our conditions. To address this, we introduce a parameter \( \theta_y \) 
and modify the definition of \( \Gamma \), where \( \theta_y \) is continuous in \( y \).
\end{remark}

\begin{remark}\label{r3} \rm
In lemma \ref{L3.17}, we prove that hypothesis (A11) is equivalent to 
$a^0\leq\inf _{x\in{\mathbb{R}^N}}[s(x)-\frac{V(x)}{g^2(\infty)}]\leq s_{\infty}-\frac{V_{\infty}}{g^2(\infty)}$, 
where $a^0>0$ is independent on the specific forms of $V(x)$ and $s(x)$.
\end{remark}

\begin{remark}\label{r4} \rm 
Assuming the decomposition $f(x,t)=s(x)f(t)$, hypotheses (A9) and (A10) are readily verified by 
direct computation. Furthermore, under the additional condition $s(x)\leq s_\infty$, 
hypothesis (A7) is also satisfied. In the special case where $g\equiv1$,  $f(t)=t-\ln(t+1)$ 
fulfills our assumptions. Note that the choice of $f(t)$ is independent of that of $s(x)$. 
As an example, for suitably chosen functions $g(t)$ and $f(t)$, consider  $V(x)\equiv V_\infty$ 
and $s(x)=s_\infty-(s_\infty-\frac{V_\infty}{g^2(\infty)}-a^0)e^{-|x|}$ 
where $s_\infty>\frac{V_\infty}{g^2(\infty)}$ and $a^0$ is as in the Lemma \ref{L3.17}. 
This choice satisfies conditions (A1) and (A11). 
\end{remark}

Prior research has predominantly focused on nonlinearities of specific types. 
For instance, reference \cite{shen} investigated the case where the nonlinear term is independent 
of the spatial variable, i.e., $f(x,t)=f(t)$, and obtained  a nontrivial Mountain Pass solution.  
Under our assumptions,  the Mountain Pass Geometry can be established by Lemmas \ref{L3.3}, 
\ref{L3.4} and the weak continuity of $J'$. However, the nontriviality of the solution cannot be 
guaranteed. To address this, we introduce additional hypotheses (A5) and (A6), which ensure that the Nehari manifold exhibits a suitable geometric structure.

Via \eqref{p3}, the problem studied in \cite{fang3} can be recovered as a particular example of 
our model, with the condition $s(x)\leq s_\infty$ there being equivalently replaced by 
$V(x)\leq V_\infty$ here, under the  substitution $g(t)\longleftrightarrow \sqrt{1+2t^2}$ and 
$f(x,t)\longleftrightarrow q(x)g(t)$ from \cite{fang3}.

\begin{remark}\label{r5} \rm
To prevent our minimax value from being an accumulation point of critical levels of 
\eqref{q4} we impose hypothesis (A12) as in \cite{Ambrosetti} (see also in \cite{Carriao}). 
\end{remark}


\noindent\textbf{Notation.} 
Throughout this paper, we use the letters \(C\), \(C_1\), \(C_2\), $\ldots$ to denote positive constants whose values may change from line to line and play no essential role. 
For a center \(y \in \mathbb{R}^N\) and a radius \(r > 0\), the open ball \(B_r(y)\) is defined as 
\(\{x \in \mathbb{R}^N : |x - y| < r\}\). The symbol \(S^*\) stands for the unit sphere in the function 
space \(E\). The usual \(L^p\)-norm is written as 
\(|u|_p := \bigl(\int_{\mathbb{R}^N} |u|^p\bigr)^{1/p}\). Moreover, given a translation vector 
\(y \in \mathbb{R}^N\) and a dilation parameter \(\theta \in \mathbb{R}^+\), we define an action of 
\(\mathbb{R}^N \times \mathbb{R}^+\) on \(E\) by setting 
\(u^{y,\theta}(x) := u\bigl(\frac{x-y}{\theta}\bigr)\).

\section{Preliminary results} \label{s2}

The variational structure associated with equation \eqref{q1} can be formally described by 
the energy functional
$$
J_0(v)=\frac{1}{2}\int_{\mathbb{R}^N}g^2(v)|\nabla v|^2+V(x)v^2-\int_{\mathbb{R}^N}F(x,v).
$$
However, $J_0$ may fail to be well‑defined   in the space $E$. We follow the approach introduced 
in \cite{shen} and perform a variable substitution,
$$
w=G(v)=\int^v_0g(\tau)\,d\tau.
$$
Hence, we obtain that
\begin{equation}
\label{q2}
J(w)=\frac{1}{2}\int_{\mathbb{R}^{N}}|\nabla w|^{2}+V(x)|G^{-1}(w)|^2-\int_{\mathbb{R}^{N}}F(x,G^{-1}(w)).
\end{equation}
From (A2), (A3) and (A4), we deduce that $J\in C^1(E,\mathbb{R})$. Obviously, we have
\begin{equation}
\label{q3}
\langle J'(w),\phi\rangle =\int_{\mathbb{R}^{N}}\nabla w \nabla \phi+V(x)
\frac{G^{-1}(w)}{g(G^{-1}(w))}\phi-\int_{\mathbb{R}^{N}}\frac{f(x,G^{-1}(w))}{g(G^{-1}(w))}\phi,
\end{equation}
for every $w, \phi\in E$. Moreover, a function \( u \in H^1(\mathbb{R}^{N}) \) is a critical point 
of \(J\) if and only if it is a weak solution of the equation
\begin{equation}\label{new}
-\Delta u+V(x)\frac{G^{-1}(u)}{g(G^{-1}(u))}=\frac{f(x,G^{-1}(u))}{g(G^{-1}(u))}, \quad u\in H^1(\mathbb{R}^{N}).
\end{equation}
It is established in \cite{shen} that \eqref{q1} and \eqref{new} are equivalent in $E$.

We introduce the Nehari manifold associated with the functional $J$ as
$$
\mathcal{N}:=\left\{u\in E\setminus \{0\}: \langle J'(u),u\rangle=0\right\},
$$ 
and we set the corresponding minimum level $c_0:=\inf_{\mathcal{N}}J$.

Next, consider the limiting problem
\begin{equation}
\label{q4}
-\Delta u= \frac{f_{\infty}(G^{-1}(u))}{g(G^{-1}(u))}-V_{\infty}\frac{G^{-1}(u)}{g(G^{-1}(u))}.
\end{equation}
 The corresponding limiting functional and Nehari manifold are defined as
\begin{gather}
\label{q5}
J_{\infty}(u)=\frac{1}{2}\int_{\mathbb{R}^{N}}|\nabla u|^{2}+V_{\infty}|G^{-1}(u)|^2-\int_{\mathbb{R}^{N}}F_{\infty}(G^{-1}(u)), \\
\label{q6}
\mathcal{N_{\infty}}:=\left\{u\in E\setminus \{0\}: \langle J_{\infty}'(u),u\rangle=0\right\}.
\end{gather}  

Equation \eqref{q4} admits a positive radial ground state solution 
$u_{\infty} \in C^2(\mathbb{R}^N)$, since the function 
$g^*(u):=\frac{f_{\infty}(G^{-1}(u))}{g(G^{-1}(u))}-V_{\infty}\frac{G^{-1}(u)}{g(G^{-1}(u))}$ can be 
verified to satisfy all assumptions of \cite[Theorem 3.1]{colin}, which  relies on the classical 
framework of \cite{Berestycki2} and \cite{Berestycki} (cf. \cite{colin}). 
Set $c_\infty:=J_\infty(u_\infty)$.


Let 
\begin{gather*}
\mathcal{T}:=\big\{u\in E: \int_{\mathbb{R}^N}|\nabla u|^2+\int_{\mathbb{R}^N}\frac{V(x)}{g^2(\infty)}u^2<\int_{\mathbb{R}^N}s(x)u^2\big\},\\
\mathcal{T_{\infty}}:=\big\{u\in E: \int_{\mathbb{R}^N}|\nabla u|^2+\int_{\mathbb{R}^N}\frac{V_\infty}{g^2(\infty)}u^2<\int_{\mathbb{R}^N}s_{\infty}u^2\big\}.
\end{gather*}
 By conditions (A5) and (A11), we have 
 $\mathcal{T} \neq \emptyset $ and $\mathcal{T}_\infty \neq \emptyset $. 
 Although we cannot guarantee that all functions in $\mathcal{T_{\infty}}$ belong to 
 $\mathcal{T}$, it is established in Lemma \ref{L3.11} that the modified functions do 
 possess this property.
 
Given that our assumptions could not ensure that $\mathcal{N}$  forms a manifold of 
class $C^1$, we consequently adopt the framework established in \cite{Szulkin1, Szulkin2}.
 

The following properties are consequences of assumption (A2). For their proofs, we refer the reader ito \cite{shen, deng}.

\begin{lemma}\label{L2.1} The following statements hold for the function $G(t)$ and $G^{-1}(t)$:
\begin{itemize}
  \item[(1)] $G(t)$ and $G^{-1}(t)$ are odd;
  \item[(2)]  $G^{-1}(t)\leq t\leq G(t)\leq g(t)t$, for any $t>0$;
  \item[(3)] $\frac{G^{-1}(t)}{t}$ is decreasing on $t\in(0,\infty)$;
  \item[(4)] $\lim_{t\to 0}\frac{G^{-1}(t)}{t}=1$, $\lim_{t\to \infty}\frac{G^{-1}(t)}{t}=\frac{1}{g(\infty)}$.
\end{itemize}
\end{lemma}

\section{Proof of Theorem \ref{thm1.1}}  \label{s3}

\begin{lemma}\label{L3.1} 
(1)  $0\leq F(x,t)\leq \frac{f(x,t)G(t)}{2g(t)}$, for all $t\in \mathbb{R}$, $x\in{\mathbb{R}^N}$.

(2) For each $\varepsilon >0$, $p\geq1$, there is $C_{\varepsilon, p}>0$ such that
$|f(x,t)|\leq \varepsilon |t|+C_{\varepsilon, p} g(t)|G(t)|^{p}$ for all $t\in \mathbb{R}$, $x\in{\mathbb{R}^N}$.
\end{lemma}

\begin{proof}
(1) It follows from (A4) and (A6) that $\frac{f(x,t)}{g(t)G(t)}\geq 0$ for every $t\neq0$. 
It follows from (A2), Lemma \ref{L2.1}-(1), (2) and the oddness of $f(x,\cdot)$ that 
$F(x,t)\geq 0$ for every $t\in \mathbb{R}$. By (A6), we obtain that
$$
F(x,t)=\int^t_0f(x,s)ds\leq \frac{f(x,t)}{2g(t)G(t)}\int^t_0\frac{dG^2(s)}{ds}ds
=\frac{f(x,t)G(t)}{2g(t)}.
$$

(2) Let $s^*:=\sup_{x\in{\mathbb{R}^N}}s(x)$, it follows from (A2), (A5) and (A6) that 
$$
|f(x,t)|\leq s^*g(t)|G(t)|\leq C_1+s^*g(t)|G(t)|^p\leq s^*(C_2+g(t)|G(t)|^p),
$$ 
for every $p\geq1$. Using (A4), we obtain  the conclusion.
\end{proof}

For each element  \( u \in \mathcal{T} \), we introduce the auxiliary function 
\( Q_u  \) by \( Q_u(t) = J(tu) \) for $t>0$.


\begin{lemma}\label{L3.2}
(1) There is a unique $t_u>0$ such that $Q'_u(t_u)=0$. Furthermore,  $t_0u\in\mathcal{N}$  is equivalent to  $t_0=t_u$.

(2) $\mathbb{R}^+\mathcal{N}=\mathcal{T}$.
\end{lemma}

\begin{proof}
 (1) From (A4), we obtain that
$\lim_{t\to 0}\frac{F(x,t)}{|G(t)|^2}=0$.
Consider that
$$
\frac{Q_u(t)}{t^2}=\frac{1}{2}\int_{\mathbb{R}^N}|\nabla u|^2+\frac{1}{2}\int_{u\neq 0}V(x)\frac{|G^{-1}(tu)|^2}{t^2u^2}u^2-\int_{u\neq 0}\frac{F(x,G^{-1}(tu))}{t^2u^2}u^2.
$$
Applying Lemma \ref{L2.1}-(4) yields 
$$
\lim_{t\to 0}\frac{Q_u(t)}{t^2}=\frac{1}{2}\int_{\mathbb{R}^N}|\nabla u|^2+V(x)u^2>0.
$$
 Using (A5), we have
$\lim_{t\to\infty}\frac{F(x,t)}{|G(t)|^2}=\frac{1}{2}s(x)$. Thus,
$$
\lim_{t\to\infty}\frac{Q_u(t)}{t^2}= \frac{1}{2}\int_{\mathbb{R}^N}|\nabla u|^2
+\frac{V(x)}{g^2(\infty)}u^2-\frac{1}{2}\int_{\mathbb{R}^N}s(x)u^2<0.
$$
Obviously, the condition $Q'_u(t) = 0$ is equivalent to
$$
\int_{\mathbb{R}^N}|\nabla u|^2+\int_{u\neq 0}V(x)\frac{G^{-1}(tu)}{g(G^{-1}(tu))tu}u^2
=\int_{u\neq 0}\frac{f(x,G^{-1}(tu))}{g(G^{-1}(tu))tu}u^2.
$$
It follows from (A2) and Lemma \ref{L2.1}-(3) that 
$$
\frac{G^{-1}(\tau)}{g(G^{-1}(\tau))\tau}=\frac{G^{-1}(\tau)}{\tau}\cdot \frac{1}{g(G^{-1}(\tau))}
$$
decreases for $\tau > 0$. Moreover, $\frac{d}{d\tau}(\frac{G^{-1}(\tau)}{\tau})|_{\tau=\tau_0}=0$ 
if and only if $g(\tau)=1$ for any $0\leq |\tau|\leq |\tau_0|$. Using (A6), we obtain that 
$\frac{f(x,G^{-1}(\tau))}{g(G^{-1}(\tau))\tau}-\frac{V(x)G^{-1}(\tau)}{g(G^{-1}(\tau))\tau}$ 
is strictly increasing on $\tau\in(0,\infty)$ and strictly decreasing on $\tau\in(-\infty,0)$, 
which yields the existence and uniqueness of $t_u$. The conclusion follows directly from  the r
elation $Q'_u(t) = \frac{\langle J'(tu), tu \rangle}{t}$.
 
(2) Obviously, $\mathcal{T}\subset\mathbb{R}^+\mathcal{N}$. Remember that  
$\frac{f(x,G^{-1}(\tau))}{g(G^{-1}(\tau))\tau}-\frac{V(x)G^{-1}(\tau)}{g(G^{-1}(\tau))\tau}$ is strictly 
increasing on $\tau\in(0,\infty)$ and strictly decreasing on $\tau\in(-\infty,0)$, for each  \( u \in \mathcal{N} \), we obtain
\begin{align*}
\int_{\mathbb{R}^N}|\nabla u|^2 
&= \int_{u\neq 0}\frac{f(x,G^{-1}(u))}{g(G^{-1}(u))u}u^2-V(x)\frac{G^{-1}(u)}{g(G^{-1}(u))u}u^2\\
&< \int_{\mathbb{R}^N}s(x)u^2-\frac{V(x)}{g^2(\infty)}u^2.
\end{align*}
Hence $u\in\mathcal{T}$, which means $\mathcal{N}\subset\mathcal{T}$, i.e., 
$\mathbb{R}^+\mathcal{N}\subset\mathcal{T}$.
\end{proof}


\begin{lemma}\label{L3.3} 
(1) There exist positive constants $r_0$ and $M$ such that $M\leq \inf_{S_{r_0}}J\leq c_0,$ where $S_{r_0}:=\{u\in E:\|u\|=r_0\}$.

(2) Every $u\in \mathcal{N}$ satisfies $\|u\|^{2}\geq 2c_0$.
\end{lemma}

\begin{proof}
(1) The estimate $\inf_{S_{r_0}}J\geq M$ for some $M>0$ follows immediately from 
\cite[Lemma 2.1]{shen},  so its proof is omitted. Using Lemma \ref{L3.2}, for every 
$u\in\mathcal{N}$, we have $\frac{r_0u}{\|u\|}\in S_{r_0}$ and 
$J(u)\geq J(\frac{r_0u}{\|u\|})$. Consequently, the first inequality holds.

(2)  The conclusion is obtained directly from $|G^{-1}(u)|^2\leq u^2$ and $F(x,t)\geq0$.
\end{proof}


\begin{lemma}\label{L3.4} 
For every compact subset  $\mathcal{B}\subset \mathcal{T}$, there is $r_1>0$ such that 
$J(tu)\leq 0$ for any $u\in\mathcal{B}$ and $t\geq r_1$.
\end{lemma}

\begin{proof}
Arguing by contradiction, set $(u_n)\subset \mathcal{B}$ and $(t_n)\subset \mathbb{R}^+$ with 
$J(t_nu_n)\geq 0$ and  $t_n\to\infty$. In addition, because of the compactness of 
$\mathcal{B}$, there exists $u \in \mathcal{B}$ such that, up to a subsequence, 
$u_n$ converges to $u$. Then, applying (A5), Lemma \ref{L2.1}-(4)  and the Lebesgue 
dominated convergence theorem, we obtain
\begin{align*}
0\leq\frac{J(t_nu_n)}{t^2_n}
&=\frac{1}{2}\int_{\mathbb{R}^N}|\nabla u_n|^2
+\frac{1}{2}\int_{u_n\neq 0}V(x)\frac{|G^{-1}(t_nu_n)|^2}{t^2_nu^2_n}u^2_n
 -\int_{u_n\neq 0}\frac{F(x,G^{-1}(t_nu_n))}{t^2_nu^2_n}u^2_n \\
&\to \frac{1}{2}\int_{\mathbb{R}^N}|\nabla u|^2+\frac{V(x)}{g^2(\infty)}u^2-\frac{1}{2}\int_{\mathbb{R}^N}s(x)u^2<0,
\end{align*}
a contradiction.
\end{proof}

The preceding lemmas hold analogously for the limiting functional $J_\infty$, 
with the replacements of $\mathcal{N}$, $\mathcal{T}$, and $c_0$ by $\mathcal{N_\infty}$, 
$\mathcal{T_\infty}$, and $c_\infty$, respectively.

\begin{lemma}\label{L3.5} 
In the Nehari manifold $\mathcal{N}$, every Palais‑Smale sequence $(v_n)$  is bounded.
\end{lemma}

\begin{proof}
Argument by contradiction. Suppose that there exists a Palais–Smale sequence \((v_n)\) 
contained in \(\mathcal{N}\) such that \(J(v_n)\le C\) and \(J'(v_n)\to0\) while \(\|v_n\|\to\infty\). 
Define the normalized sequence \(w_n := v_n/\|v_n\|\). By extracting a suitable subsequence, 
\(w_n\) converges weakly in \(E\) to  \(w\), and converges almost everywhere in \(\mathbb{R}^N\) t
o $w$.

If
$$
\lim_{n\to\infty}\sup_{y\in{\mathbb{R}^N}}\int_{B_1(y)}w^2_n\to 0,
$$
then according to P.L. Lions' lemma \cite[p.16]{Willem}, 
 we obtain $w_n\to 0$ in $L^q({\mathbb{R}^N})$ for any $q\in (2, 2^*)$, with $2^*:= \frac{2N}{N-2}$ for 
 $N \geq 3$ and $2^* := \infty$ otherwise. 
 From Lemma \ref{L2.1}-(2) and Lemma \ref{L3.1}-(2),  for every fixed  $t>0$, we find 
 $\int_{\mathbb{R}^N}F(x,G^{-1}(t w_n))\to 0$ as $n\to\infty$. Using Lemma \ref{L2.1}-(4), we have
 $$
 \lim_{\tau\to 0}\frac{\tau^2-|G^{-1}(\tau)|^2}{\tau^2}=0\quad \text{and}\quad
 \lim_{\tau\to \infty}\frac{\tau^2-|G^{-1}(\tau)|^2}{\tau^p}=0.
 $$
Thus, for every $\varepsilon>0$, there exists $C_{\varepsilon}>0$ such that 
$\left|t^2w^2_n-|G^{-1}(t^2w^2_n)|^2\right|\leq \varepsilon |tw_n|^2+C_{\varepsilon}|tw_n|^p$. 
Hence $$\int_{\mathbb{R}^N}V(x)\left(t^2w^2_n-|G^{-1}(tw_n)|^2\right)\to 0,$$ for every $t>0$. 
Using Lemma \ref{L3.2}, we obtain $J(tw_n)\leq J(v_n)\leq d$ and
\begin{align*}
 J(tw_n)&= \frac{t^2}{2}\int_{\mathbb{R}^{N}}|\nabla w_n|^{2}
  +\frac{1}{2}\int_{\mathbb{R}^N}V(x)|G^{-1}(tw_n)|^2-\int_{\mathbb{R}^{N}}F(x,G^{-1}(tw_n)) \\
 &= \frac{t^2}{2}\|w_n\|^2-\frac{1}{2}\int_{\mathbb{R}^N}V(x)\left(t^2w^2_n-|G^{-1}(tw_n)|^2\right)
  -\int_{\mathbb{R}^{N}}F(x,G^{-1}(tw_n)).
\end{align*}
For sufficiently large $t$, this leads to a contradiction.

Consequently, taking a subsequence if necessary, we can find a constant $\eta > 0$ and a 
sequence $(y_n) \subset \mathbb{R}^N$ satisfying
$$
\int_{B_1(0)}(w_n^{y_n,1})^2\geq \eta>0.
$$
Provided that the sequence $(y_n)$ is bounded, then  $y_n\to y_0$ by extracting a suitable 
subsequence which informs $w_n^{y_n,1}\to w^{y_0,1}$ in $L^2(B_1(0))$ and 
$w(x) \not\equiv 0$. Furthermore, since $v_n(x) \to \infty$ on the set $\{x: w(x) \neq 0\}$, 
for every $\varphi \in C_0^{\infty}(\mathbb{R}^N)$, it holds that 
$\left\langle \frac{J'(v_n)}{\|v_n\|}, \varphi \right\rangle \to 0$, i.e.,
$$
\int_{\mathbb{R}^N}\nabla w_n\nabla\varphi+\int_{v_n
\neq 0}\frac{V(x)G^{-1}(v_n)}{g(G^{-1}(v_n))v_n}w_n\varphi
=\int_{v_n\neq 0}\frac{f(x,G^{-1}(v_n))}{g(G^{-1}(v_n))v_n}w_n\varphi+o(1),
$$
where $o(1)\to 0$ as $n\to\infty$.
Since \(w_n\) converges weakly in \(E\) to  \(w\), we obtain
$$
\int_{\mathbb{R}^N}\nabla w\nabla\varphi+\frac{V(x)}{g^2(\infty)}w\varphi=\int_{\mathbb{R}^N}s(x)w\varphi.
$$
Considering that the essential spectrum of $-\Delta-\left(s(x)-\frac{V(x)}{g^2(\infty)}\right)$ is 
$[-s_{\infty}+\frac{V_{\infty}}{g^2(\infty)},\infty)$ (cf. Theorem 3.15 in \cite{Stuart}),  
the preceding result leads to a contradiction.

Then $|y_n|\to\infty$,  it follows that, up to a subsequence, 
$w_n^{y_n,1} \rightharpoonup w^* \not\equiv 0$. Likewise, we have  
$$
\int_{\mathbb{R}^N}\nabla w^*\nabla\varphi+\frac{V_\infty}{g^2(\infty)}w^*\varphi=\int_{\mathbb{R}^N}s_\infty w^*\varphi,
$$ 
and it is also a contradiction. The conclusion follows.
\end{proof}

\begin{lemma}\label{L3.6} 
If $g\not\equiv1$, then $\lim_{t\to\infty}t^2-\frac{G^2(t)}{g^2(\infty)}=\infty$.
\end{lemma}

\begin{proof}
We  shall only  consider $t>0$. By (A2), we have 
$$
t^2-\frac{G^2(t)}{g^2(\infty)}\geq t^2-\frac{tG(t)}{g(\infty)}
=t \int_{0}^{t}[1-\frac{g(s)}{g(\infty)}]ds.
$$
If  $g\not\equiv1$, then there is a positive constant $C$ such that 
$t^2-\frac{G^2(t)}{g^2(\infty)}\geq Ct$ for any $t>1$. The conclusion follows.
\end{proof}


We define $\mathcal{N}_0:=\mathcal{T}\cap S^*$. We introduce a map 
$n:\mathcal{N}_0\mapsto \mathcal{N}$ defined by $n(u):=t_uu$, where $t_u$ is given 
by Lemma \ref{L3.2}-(1). The openness of $\mathcal{T}$ in $E$ implies the openness of 
$\mathcal{N}_0$ in $S^*$.

\begin{lemma}\label{L3.7} 
If $(u_n)\subset \mathcal{N}_0$, $t_nu_n\in\mathcal{N}$ and 
$u_n\to u_0\in \partial \mathcal{N}_0$, then $t_n\to \infty$ and $J(t_nu_n)\to \infty$.
\end{lemma}

\begin{proof}
Suppose $t_n\not\to\infty$, then we can assume $t_n\to t_0\geq 2 c_0>0$ from 
Lemma \ref{L3.3}-(2).
It follows from $0=\langle J'(t_nu_n),t_nu_n\rangle\to \langle J'(t_0u_0),t_0u_0\rangle$ 
that  $t_0u_0\in\mathcal{N}$, which is contrary to $u_0\in \partial \mathcal{N}_0$.

Note that
$$
\int_{\mathbb{R}^N}|\nabla u_0|^2+\int_{\mathbb{R}^N}\frac{V(x)}{g^2(\infty)}u^2_0=\int_{\mathbb{R}^N}s(x)u^2_0.
$$
Thus, we have 
\begin{align*}
 J(tu_0)&= \frac{t^2}{2}\int_{\mathbb{R}^{N}}|\nabla u_0|^{2}+\frac{1}{2}\int_{\mathbb{R}^N}V(x)|G^{-1}(tu_0)|^2-\int_{\mathbb{R}^{N}}F(x,G^{-1}(tu_0)) \\
        &= \frac{1}{2}\int_{\mathbb{R}^N}V(x)\left[|G^{-1}(tu_0)|^2-\frac{t^2u^2_0}{g^2(\infty)}\right]
        +\int_{\mathbb{R}^{N}}\frac{1}{2}s(x)t^2u^2_0-F(x,G^{-1}(tu_0)).
\end{align*}
Using (A8), Lemma \ref{L3.6} and Fatou's lemma, we have $J(tu_0)\to \infty$, as $t\to\infty$. 
For any positive integer $k$, choose $t_k>0$ that satisfies $J(t_ku_0)\geq k$. 
From Lemma \ref{L3.2}-(1), we have $J(t_nu_n)\geq J(t_ku_n)\to J(t_ku_0)\geq k$. 
Since $k$ was arbitrary, it follows that $J(t_nu_n)\to \infty$, as $n\to\infty$.
\end{proof}


While the functional $J$ is not coercive on $\mathcal{N}$ under our assumptions 
(whereas coercivity would imply the continuity of $n$ by \cite[Lemma 2.8]{Szulkin1}), the following lemma states that $n$ is continuous.


\begin{lemma}\label{L3.8} 
The map $n:\mathcal{N}_0\mapsto \mathcal{N}$ is continuous.
\end{lemma}

\begin{proof}
Set $(u_n)\subset \mathcal{N}_0$  satisfy $u_n\to u_0$ and let $n(u_n)=t_nu_n$. 
The sequence $(t_n)$ is bounded. Indeed, if it were unbounded, we would have 
(after passing to a subsequence) $t_n \to \infty$. It follows from  (A5), 
Lemma \ref{L2.1}-(4) and the Lebesgue dominated convergence theorem that
\begin{align*}
0&\leq\frac{J(t_nu_n)}{t^2_n}\\
&= \frac{1}{2}\int_{\mathbb{R}^N}|\nabla u_n|^2+\frac{1}{2}\int_{u_n\neq 0}V(x)\frac{|G^{-1}(t_nu_n)|^2}{t^2_nu^2_n}u^2_n
-\int_{u_n\neq 0}\frac{F(x,G^{-1}(t_nu_n))}{t^2_nu^2_n}u^2_n \\
&\to  \frac{1}{2}\int_{\mathbb{R}^N}|\nabla u_0|^2+\frac{V(x)}{g^2(\infty)}u_0^2-\frac{1}{2}\int_{\mathbb{R}^N}s(x)u_0^2<0.
\end{align*}
Thus, along a subsequence, we obtain $t_n\to t^*$. From Lemma \ref{L3.3}-(2), we 
have $t^*\geq 2c_0>0$. Hence 
$$
0=\langle J'(t_nu_n),t_nu_n\rangle\to \langle J'(t^*u_0),t^*u_0\rangle,
$$  
which implies $t^*u_0=n(u_0)$.
\end{proof}

\begin{lemma}\label{L3.9} 
The map $n:\mathcal{N}_0\mapsto \mathcal{N}$ is a homeomorphism whose inverse  is  
$n^{-1}(u)=\frac{u}{\|u\|}$.
\end{lemma}

\begin{proof}
The invertibility of $n$  follows directly from Lemma \ref{L3.2}.
Using Lemma \ref{L3.3}-(2), we have 
$$
\|n^{-1}(w_1)-n^{-1}(w_2)\|=\|\frac{w_1-w_2}{\|w_1\|}+\frac{(\|w_2\|-\|w_1\|)w_2}{\|w_1\|\|w_2\|}\|\leq \sqrt{ \frac{2}{c_0}}\|w_1-w_2\|,
$$
for every $w_1, w_2\in\mathcal{N}$. Then $n^{-1}$ is Lipschitz continuous.
\end{proof}

To study the functional $J$ on the Nehari manifold $\mathcal{N}$, we define 
$\Upsilon: \mathcal{N}_0\to \mathbb{R}$ by $\Upsilon(w) := J(n(w))$, 
which leads to the following lemma on the correspondence between the critical values 
and Palais-Smale sequences of $\Upsilon$ and $J$. 
A key property of $\Upsilon$ and $J$ is given in the next lemma 
(cf. \cite[Lemma 3.9]{fang3}).

\begin{lemma}\label{L3.10} 
(1) $\Upsilon\in C^{1}(\mathcal{N}_0,\mathbb{R})$ and
$$
\langle \Upsilon'(u),v \rangle =\|n(u)\| \langle J'(n(u)),v \rangle \quad \text{for all } v\in T_{u}(\mathcal{N}_0).
$$

(2) If $(v_{n})\subset\mathcal{N}_0$ is a Palais-Smale sequence for $\Upsilon$,
 then $(n(v_{n}))\subset\mathcal{N}$ is a Palais-Smale sequence for $J$. 
On the other hand, if $(u_{n})\subset\mathcal{N}$ is a bounded Palais-Smale sequence for $J$, 
then $(n^{-1}(u_{n}))\subset\mathcal{N}_0$ is a Palais-Smale sequence for $\Upsilon$.

(3) The critical points of $\Upsilon$ and the nontrivial critical points of $J$ are in a one‑to‑one 
correspondence via the map $n$: v is a critical point of $\Upsilon$ if and only if $n(v)$ is a 
nontrivial critical point of $J$. Furthermore, the corresponding critical values coincide, and we
 have $\inf_{\mathcal{N}_0}\Upsilon=\inf_{\mathcal{N}}J$.
\end{lemma}

\begin{lemma}\label{L3.11} 
$c_0\leq c_\infty$.
\end{lemma}

\begin{proof} 
Note that $u_{\infty}^{y,1}(x):=u_{\infty}(x-y)$. Using Lemma \ref{L2.1}-(4), (A3), (A5) and an 
integral Transform, we have
\begin{align*}
\frac{\langle J'_{\infty}(ru_{\infty}^{y,1}),ru_{\infty}^{y,1}\rangle}{r^{2}} 
&= \int_{\mathbb{R}^{N}} \vert \nabla u_{\infty}^{y,1} \vert^{2} + \int_{\mathbb{R}^{N}} V_{\infty} \frac{G^{-1}(ru_{\infty}^{y,1})}{g(G^{-1}(ru_{\infty}^{y,1}))ru_{\infty}^{y,1}} (u_{\infty}^{y,1})^{2} \\
&\quad - \int_{\mathbb{R}^{N}}  \frac{f_\infty (G^{-1}(ru_{\infty}^{y,1}))}{g(G^{-1}(ru_{\infty}^{y,1}))ru_{\infty}^{y,1}} (u_{\infty}^{y,1})^{2} \\
&\to \int_{\mathbb{R}^{N}} \vert \nabla u_{\infty} \vert^{2} +\int_{\mathbb{R}^{N}} \frac{V_{\infty} }{g^2(\infty)}u_{\infty}^{2}- \int_{\mathbb{R}^{N}} s_{\infty} u_{\infty}^{2} < 0,
\end{align*}
as $r\to\infty$, for every $y\in{\mathbb{R}^N}$.  Consequently, for some $\alpha<0$ and $R\geq0$, 
the inequality $\frac{\langle J'_{\infty}(ru_{\infty}^{y,1}),ru_{\infty}^{y,1}\rangle}{r^{2}}\leq \alpha$ 
holds for any $r\geq R$.  Applying  (A1) and Lemma \ref{L2.1}-(3), (4), we obtain
$$
\lim_{|y|\to\infty} \int_{\mathbb{R}^{N}} [V(x+y)-V_{\infty}]\frac{G^{-1}(ru_{\infty})}{g(G^{-1}(ru_{\infty}))ru_{\infty}} u_{\infty}^{2}=0.
$$
A similar argument shows that 
$$
\lim_{|y|\to\infty}  \int_{\mathbb{R}^{N}}  \frac{f_\infty (G^{-1}(ru_{\infty}))-f(x+y,G^{-1}(ru_{\infty}))}{g(G^{-1}(ru_{\infty}))ru_{\infty}} u_{\infty}^{2} =0,
$$
by (A3), (A5) and  (A6). 
Then $\frac{\langle J'(ru_{\infty}^{y,1}),ru_{\infty}^{y,1}\rangle}{r^{2}}
=\frac{\langle J'_{\infty}(ru_{\infty}^{y,1}),ru_{\infty}^{y,1}\rangle}{r^{2}}+o(1)$, 
where $o(1)\to 0$, as $|y|\to\infty$. Hence, for some $S>0$, the inequality 
$\frac{\langle J'(ru_{\infty}^{y,1}),ru_{\infty}^{y,1}\rangle}{r^{2}}\leq\frac{\alpha}{2}<0$ 
holds whenever $r\geq R$ and $|y|\geq S$.

Observe that $\langle J'(ru_{\infty}^{y,1}),ru_{\infty}^{y,1}\rangle$ is positive definite for 
small $r>0$. Following the argument in Lemma \ref{L3.2}-(1), we find that 
$\frac{\langle J'(ru_{\infty}^{y,1}),ru_{\infty}^{y,1}\rangle}{r^{2}}$ decreases strictly on 
$r\in(0,\infty)$.
Consequently, for every $|y|\geq S$, the interval $(0, R)$ contains precisely one element 
$T^y$ satisfying   $\langle J'(T^yu_{\infty}^{y,1}),T^yu_{\infty}^{y,1}\rangle=0$, which means 
$T^y u_{\infty}^{y,1}\in\mathcal{N}$. Thus,
\begin{align*}
\int_{\mathbb{R}^{N}} \vert \nabla u_{\infty} \vert^{2} 
&= \int_{\mathbb{R}^{N}} [\frac{f(x+y,G^{-1}(T^{y}u_{\infty}))}{g(G^{-1}(T^{y}u_{\infty}))
 T^{y}u_{\infty}} - \frac{V(x + y)G^{-1}(T^{y}u_{\infty})}{g(G^{-1}(T^{y}u_{\infty}))T^{y}u_{\infty}}]u_{\infty}^{2} \\
&= \int_{\mathbb{R}^{N}} [\frac{f_\infty(G^{-1}(T^{y}u_{\infty}))}{g(G^{-1}(T^{y}u_{\infty}))T^{y}u_{\infty}} - \frac{V_\infty G^{-1}(T^{y}u_{\infty})}{g(G^{-1}(T^{y}u_{\infty}))T^{y}
u_{\infty}}]u_{\infty}^{2} + o(1).
\end{align*} 
Given that $\frac{f_\infty(G^{-1}(s))}{g(G^{-1}(s))s} -\frac{V_\infty G^{-1}(s)}{g(G^{-1}(s))s}$  
is strictly increasing on $s\in(0,\infty)$, we have $T^y\to1$ and then 
$c_0\leq J(T^y u_{\infty}^{y,1})\to c_\infty$.
\end{proof}

\begin{lemma}\label{L3.12}
For each Palais-Smale sequence  $(u_n) \subset \mathcal{N}$ with $J(u_n)\to C$, 
after passing to a suitable subsequence, there exists a solution $u^0 \in E$ to
\[
-\Delta w + V(x)\frac{G^{-1}(w)}{g(G^{-1}(w))}= \frac{f(x,G^{-1}(w))}{g(G^{-1}(w))},
\]
 a finite sequence of solutions $u^1, \ldots, u^k \in E$ to
\[
-\Delta w + V_\infty \frac{G^{-1}(w)}{g(G^{-1}(w))} = \frac{f_\infty(G^{-1}(w))}{g(G^{-1}(w))},
\]
and $k$ sequences $(y_n^j) \subset \mathbb{R}^N$ satisfying
\begin{gather*}
|y_n^j| \to \infty, \quad |y_n^j - y_n^{j'}| \to \infty, \quad j \neq j', \quad n \to \infty, \\
\int_{\mathbb{R}^N} |\nabla (u_n - u^0 - \sum_{j=1}^k u^j(\cdot - y_n^j))|^2 + |G^{-1}(u_n - u^0 - \sum_{j=1}^k u^j(\cdot - y_n^j))|^2\to 0, \\
\int_{\mathbb{R}^N} |\nabla u_n|^2 + |G^{-1}(u_n) |^2
\to \sum_{j=0}^k \int_{\mathbb{R}^N} |\nabla u^j|^2 + |G^{-1}(u^j)|^2, \\
J(u^0) + \sum_{j=1}^k J_\infty (u^j) = C.
\end{gather*}
\end{lemma}

\begin{proof}  
Using Lemma \ref{L3.5},  selecting a subsequence if necessary,  $u_n\rightharpoonup u^0$ in $E$. 
 We claim that $\int_{\mathbb{R}^N}|G^{-1}(u_n)|^2-|G^{-1}(u_n-u^0)|^2-|G^{-1}(u^0)|^2=o(1)$ and 
 $\int_{\mathbb{R}^N}[\frac{G^{-1}(u_n)}{g(G^{-1}(u_n))}-\frac{G^{-1}(u_n-u^0)}{g(G^{-1}(u_n-u^0))}
 -\frac{G^{-1}(u^0)}{g(G^{-1}(u^0))}]\varphi=o(1)$ where $o(1)\to0$, as $n\to\infty$, 
 uniformly for $\|\varphi\|\leq1$.

We now establish the second assertion. Since $\frac{G^{-1}(t)}{g(G^{-1}(t))t}$ is nonincreasing 
on $t\in(0,\infty)$, we have $[\frac{G^{-1}(t)}{g(G^{-1}(t))}]'\leq \frac{G^{-1}(t)}{g(G^{-1}(t))t}\leq1$, 
for any $t>0$. An application of the mean value theorem and H\"older's inequality shows that 
$$
\int_{|x|>r}[\frac{G^{-1}(u_n)}{g(G^{-1}(u_n))}-\frac{G^{-1}(u_n-u^0)}{g(G^{-1}(u_n-u^0))}]
\varphi\leq|\varphi|_{2}[\int_{|x|>r}|u^0|^2]^\frac{1}{2}.
$$ 
Thus, for each $\varepsilon > 0$, there is a corresponding $r_1 > 0$ satisfying 
$$
\int_{|x|>r_1}[\frac{G^{-1}(u_n)}{g(G^{-1}(u_n))}-\frac{G^{-1}(u_n-u^0)}{g(G^{-1}(u_n-u^0))}-\frac{G^{-1}(u^0)}{g(G^{-1}(u^0))}]\varphi\leq\varepsilon\|\varphi\|.
$$  
Then, using the Rellich theorem, we obtain that 
$$
\int_{|x|\leq r_1}[\frac{G^{-1}(u_n)}{g(G^{-1}(u_n))}-\frac{G^{-1}(u_n-u^0)}{g(G^{-1}(u_n-u^0))}-\frac{G^{-1}(u^0)}{g(G^{-1}(u^0))}]\varphi\leq\varepsilon\|\varphi\|.
$$  
The first statement can be proved similarly.

Following the argument in  Lemmas 3.4 and 3.5 of \cite{Yang}, we obtain, by (A2) and the 
proof of Lemma \ref{L3.1}-(2),  that
$\int_{\mathbb{R}^N}F(x,G^{-1}(u_n))-F(x,G^{-1}(u_n-u^0))-F(x,G^{-1}(u^0))=o(1)$ and
\[
\int_{\mathbb{R}^N}[\frac{f(x,G^{-1}(u_n))}{g(G^{-1}(u_n))}
-\frac{f(x,G^{-1}(u_n-u^0))}{g(G^{-1}(u_n-u^0))}-\frac{f(x,G^{-1}(u^0))}{g(G^{-1}(u^0))}]\varphi=o(1), 
\]
where $o(1)\to0$, as $n\to\infty$, uniformly for $\|\varphi\|\leq1$. 
Hence, $J(u_n-u^0)=J(u_n)-J(u^0)+o(1)$ and $J'(u_n-u^0)=J'(u_n)-J'(u^0)+o(1)$. 
Similarly, we obtain $J_\infty(u_n-u^0)=J_\infty(u_n)-J_\infty(u^0)+o(1)$  and 
$J'_\infty(u_n-u^0)=J'_\infty(u_n)-J'_\infty(u^0)+o(1)$. 
 
Let $u^1_n:=u_n-u^0$. Then $u^1_n\to 0$ in $L^2_{loc}({\mathbb{R}^N})$ by the Rellich theorem. 
We assert that  $J_\infty(u^1_n)=J(u_n)-J(u^0)+o(1)$ and $J'_\infty(u_n-u^0)=J'(u_n)-J'(u^0)+o(1)$.

We now prove the second assertion.
 \begin{align*}
|\langle J'_\infty(u^1_n)-J'(u^1_n),\varphi\rangle|
&\leq \int_{\mathbb{R}^N}|V(x)-V_\infty|\cdot |\frac{G^{-1}(u^1_n)}{g(G^{-1}(u^1_n))}|\cdot |\varphi|\\
&\quad +\int_{\mathbb{R}^N}|\frac{f(x,G^{-1}(u^1_n))}{g(G^{-1}(u^1_n))}-\frac{f_\infty(G^{-1}(u^1_n))}{g(G^{-1}(u^1_n))}|\cdot |\varphi|.
 \end{align*}
For each \(\varepsilon > 0\), a positive number \(r_2\) can be found such that 
\(|V(x) - V_\infty| < \varepsilon\) whenever \(|x| > r_2\). Employing H\"older's inequality under 
this condition produces 
$$
\int_{|x|>r_2}|V(x)-V_\infty|\cdot |\frac{G^{-1}(u^1_n)}{g(G^{-1}(u^1_n))}|\cdot |\varphi|\leq\varepsilon C|\varphi|_2.
$$
From the Rellich theorem, we obtain
$$
\int_{|x|\leq r_2}|V(x)-V_\infty|\cdot |\frac{G^{-1}(u^1_n)}{g(G^{-1}(u^1_n))}|\cdot |\varphi|\leq\varepsilon \|\varphi\|,
$$ 
for $n$ large enough.
Note that 
\begin{align*}
&\int_{\mathbb{R}^N}|\frac{f(x,G^{-1}(u^1_n))}{g(G^{-1}(u^1_n))}
 -\frac{f_\infty(G^{-1}(u^1_n))}{g(G^{-1}(u^1_n))}|\cdot |\varphi|\\
&\leq\int_{|x-y|\leq r_3}|\frac{f(x,G^{-1}(u^1_n))}{g(G^{-1}(u^1_n))}|\cdot
  |\varphi|+|\frac{f_\infty(G^{-1}(u^1_n))}{g(G^{-1}(u^1_n))}|\cdot |\varphi|\\
&\quad+\int_{|x-y|>r_3,\,|G^{-1}(u^1_n)|>T}|\frac{f(x,G^{-1}(u^1_n))}{g(G^{-1}(u^1_n))u^1_n}-s(x)|\cdot |u^1_n|\cdot |\varphi|\\
&\quad+\int_{|x-y|>r_3,\,|G^{-1}(u^1_n)|>T}|\frac{f_\infty(G^{-1}(u^1_n))}{g(G^{-1}(u^1_n))u^1_n}-s_\infty|\cdot |u^1_n|\cdot |\varphi|+|s_\infty-s(x)|\cdot |\varphi|\\
&\quad+\int_{|x-y|>r_3,\,|u^1_n|\leq G(T)}|\frac{f(x,G^{-1}(u^1_n))}{g(G^{-1}(u^1_n))}-\frac{f_\infty(G^{-1}(u^1_n))}{g(G^{-1}(u^1_n))}|\cdot |u^1_n|\cdot |\varphi|.
\end{align*}
Using (A3), (A5), (A9) (or (A10)), the Rellich theorem and H\"older inequality, there exist  some $T>0$ large enough (if (A9) holds, $y=0$; if (A10) holds, $y$ is as in (A10)) and $r_3>0$ large enough such that $|\langle J'_\infty(u^1_n)-J'(u^1_n),\varphi\rangle|\leq\varepsilon C \|\varphi\|$, 
for $n$ large enough. A similar argument applies to the first statement.

Set
$$
\delta:=\lim_{n\to\infty}\sup_{y\in{\mathbb{R}^N}}\int_{B_1(y)}|u^1_n|^2.
$$
In the case $\delta = 0$, we  have \(u^1_n \to 0\) in \(L^q(\mathbb{R}^N)\) for all 
\(q \in (2, 2^*)\), by the P.L. Lions lemma \cite[Lemma 1.21 ]{Willem}). 
This convergence, combined with assumptions \((h_1)\), \((g)\) and the results of 
Lemma \ref{L2.1}-(2) and Lemma \ref{L3.1}-(2), implies that for any \(\varepsilon >0\) 
one can find \(C_{\varepsilon}>0\) such that 
\begin{gather*}
\int_{\mathbb{R}^N}|\frac{f(x,G^{-1}(u^1_n))}{g(G^{-1}(u^1_n))}u^1_n|\leq\int_{\mathbb{R}^N}\varepsilon|u^1_n|^2+C_\varepsilon|u^1_n|^p,
\\
\int_{\mathbb{R}^N}|\frac{f_\infty(G^{-1}(u^1_n))}{g(G^{-1}(u^1_n))}u^1_n|\leq\int_{\mathbb{R}^N}\varepsilon|u^1_n|^2+C_\varepsilon|u^1_n|^p.
\end{gather*}
From the identity $\langle J'_\infty(u^1_n),u^1_n\rangle=\langle J'(u_n),u^1_n\rangle+\langle J'(u^0),u^1_n\rangle+o(1)=o(1)$, 
together with the estimate provided in Lemma \ref{L2.1}-(2), we obtain that
\begin{align*}
&\frac{1}{2}\int_{\mathbb{R}^N}|\nabla u^1_n|^2+V_\infty|G^{-1}(u^1_n)|^2 \\
&\leq \frac{1}{2}\int_{\mathbb{R}^N}|\nabla u^1_n|^2+\frac{1}{2}\int_{\mathbb{R}^N}V_\infty|G^{-1}(u^1_n)|^p
 +C\int_{|G^{-1}(u^1_n)|\leq1}V_\infty\frac{G^{-1}(u^1_n)u^1_n}{g(G^{-1}(u^1_n))}\\
&\leq C\int_{\mathbb{R}^N}|\nabla u^1_n|^2+V_\infty\frac{G^{-1}(u^1_n)u^1_n}{g(G^{-1}(u^1_n))}
 +\frac{V_\infty}{2}\int_{\mathbb{R}^N}|u^1_n|^p \to  0,
\end{align*}
as $n\to\infty$. The conclusion follows.

If $\delta>0$, then there is a sequence \((y^1_n)\) in \(\mathbb{R}^N\) for which 
$\int_{B_1(0)}|(u^1_n)^{-y_n,1}|^2>\frac{\delta}{2}$. By the local compactness of the 
Sobolev embedding and after passing to a subsequence, we may assume that 
$(u^1_n)^{-y_n,1} \rightharpoonup u^1\neq0$. The weak convergence 
$u^1_n \rightharpoonup 0$ implies that $(y^1_n)$ is unbounded. Therefore, we may extract a subsequence such that $y^1_n \to \infty$, as $n \to \infty$. 
The identity $J'_\infty(u^1)=0\in H^{-1}({\mathbb{R}^N})$  holds due to the weak sequential continuity 
and translation invariance of $J'_\infty$. 

Define $u^2_n:=u^1_n-(u^1)^{y_n,1}$, then $u^2_n \rightharpoonup 0$ and 
$(u^2_n)^{-y_n,1} \rightharpoonup 0$ in $E$. Take the same argument as above and repeat 
the process, except that $J$ should be replaced by $J_\infty$.  
Since $(J(u_n))$ is bounded, the fact that every nontrivial critical point of $J_\infty$  have 
energy at least  $c_\infty$ yields that the iteration terminates finitely.
\end{proof}

\begin{lemma}\label{L3.13} 
For each $b>0$,  there exist a constant $C_b \in(0,1]$ such that 
$$
C_b\|u\|^2\leq \int_{\mathbb{R}^N}|\nabla u|^2+V(x)|G^{-1}(u)|^2\leq\|u\|^2,
$$ 
for all $\|u\|\leq b$.
\end{lemma}

\begin{proof}
The detailed argument can be found in \cite[Lemma 3.2 ]{fang1}. 
Nevertheless, we provide a proof sketch for the sake of completeness. 
An application of Lemma \ref{L2.1}(2) yields that 
$$
\int_{\mathbb{R}^N}|\nabla u|^2+V(x)|G^{-1}(u)|^2\leq\|u\|^2.
$$
Suppose, for a contradiction, that there is a sequence $(u_n)$ of nonzero elements in $E$ 
satisfying $\|u_n\|\leq b$ and
\[
\int_{\mathbb{R}^N} |\nabla w_n|^2 + \int_{u_n \neq 0} V(x) \frac{|G^{-1}(u_n)|^2}{u_n^2} w_n^2 \to 0,
\]
where $w_n := \frac{u_n}{\|u_n\|}$. This leads to $\int_{\mathbb{R}^N} |\nabla w_n|^2 \to 0$, whereas
\begin{equation}
\label{q7}
\int_{\mathbb{R}^N} V(x) w_n^2 \to 1. 
\end{equation}
For a given constant $C_1 > 0$, let $M^1_n := \{ x \in \mathbb{R}^N : |u_n(x)| \geq C_1 \}$ 
and $M^2_n$ denote its complement. 
Since $\|u_n\|\leq b$,  it follows that for an arbitrary  $\varepsilon > 0$, 
there exists $C_1=C_\varepsilon$  satisfying $|M^1_n| \leq \varepsilon$. 
Additionally, Lemma \ref{L2.1}-(3) implies that $\frac{G^{-1}(t)}{t}$  decreases on $t \in(0,\infty)$. Thus,
\[
\frac{|G^{-1}(C_1)|^2}{C_1^2} \int_{M^2_n} V(x) w_n^2 \leq \int_{M^2_n} V(x) \frac{|G^{-1}(u_n)|^2}{u_n^2} w_n^2 \to 0.
\]
For sufficiently small $\varepsilon$,  combining H\"older's inequality with assumption (A1), 
we obtain that
\[
\int_{M^1_n} V(x) w_n^2 \leq C_2\varepsilon^{(2^*-2)/2^*} \leq \frac{1}{2},
\]
which contradicts \eqref{q7}. 
\end{proof}

\begin{proposition} \label{Pro1} 
Suppose conditions  {\rm (A1)--(A8), (A11) } hold, and at least one of {\rm (A9) }and {\rm (A10) }
holds. Then, the strict inequality \( c_0 < c_\infty \) guarantees the existence of a nontrivial ground state solution to \eqref{new}.
\end{proposition}

\begin{proof}
Although Ekeland's variational principle (see \cite[Theorem 4.8.1]{Chang}) is stated for 
a complete metric space, Lemma \ref{L3.7} ensures its validity in $\mathcal{N}_0$, 
as it precludes the limiting point from reaching the boundary.
So, there exists $(w_n)\subset\mathcal{N}$ such that $\Upsilon (w_n)\to c_0$ and 
$\Upsilon' (w_n)\to 0$. Let $u_n:=n(w_n)$.  Then, by Lemmas \ref{L3.5} and \ref{L3.10}, 
the sequence $(u_n)$ constitutes a bounded Palais-Smale sequence for the functional $J$ 
at level $c_0$. Since $c_0 < c_\infty$,  Lemma \ref{L3.12}  implies 
$\int_{\mathbb{R}^N}|\nabla(u_n-u^0)|^2+|G^{-1}(u_n-u^0)|^2\to0$ and thus 
$\|u_n-u^0\|\to0$ by Lemma \ref{L3.13}. Using $J(n(w_n))\to J(u^0)$ and Lemma \ref{L3.7}, 
we obtain $w_n\to w^0 \in \mathcal{N}_0$ and $u^0=n(w^0)\in\mathcal{N}$.
\end{proof}

We now consider the case \( c_0 = c_\infty \). Following \cite{Ambrosetti, Cerami}, 
we define the barycenter map \( \beta \) for  \( u \in E\setminus\{0\} \). 
First, define \( \tau(u)(x) := \frac{1}{|B_1|} \int_{B_1(x)} |u(y)| \, dy \), which is a continuous 
function and bounded. Then, set \( \hat{u}(x) := \left[ \tau(u)(x) - \frac{1}{2} \max \tau(u) \right]^+\).
 The barycenter of \( u \) is 
 \[
\beta(u) := \frac{1}{|\hat{u}|_1} \int_{\mathbb{R}^N} x \hat{u}(x) \, dx \in \mathbb{R}^N.
 \]
Owing to the compact support of $\hat{u}$,  the map $\beta(u)$ defined above is well-defined 
and verifies the following:
\begin{itemize}
\item[(1)] $\beta$ is continuous in $ E\setminus\{0\}$.
\item[(2)] $\beta(u)=0$ whenever $u$ is a radial function.
\item[(3)] $\beta(t u) = \beta(u)$, for any $t > 0$.
\item[(4)] $\beta(u^{y,1})=\beta(u)+y$, for any  $y\in{\mathbb{R}^N}$.
\end{itemize}
Set $b:=\inf_{u\in\mathcal{N},\beta(u)=0}J(u)=\inf_{v\in \mathcal{N}_0,\beta(n(v))=0}\Upsilon (v)$. Directly from the definitions, it follows that  $b\geq c_0=c_\infty$.

\begin{proposition} \label{Pro2} 
Suppose conditions  {\rm (A1)--(A8), (A11) } hold, and at least one of {\rm (A9) }and {\rm (A10) }
 holds. Then,  the equality $b=c_0=c_\infty$ implies the existence of a nontrivial ground state solution to \eqref{new}.
\end{proposition}

\begin{proof}
Consider a minimizing sequence $(v_n)\subset \mathcal{N}_0$  for $\Upsilon$ satisfying 
$\beta(n(v_n))=0$ and $\Upsilon(v_n)\to b$. An application of Ekeland’s variational principle 
gives another sequence  $(w_n)\subset \mathcal{N}_0$ with $\Upsilon(w_n)\to b$, 
$\Upsilon'(w_n)\to 0$ and $\|w_n-v_n\|\to0$. Define $u_n:=n(w_n)=t_nw_n$, then 
$(u_n)\subset\mathcal{N}$ constitutes a Palais-Smale sequence for $J$ at the level $b$. 
It follows from Lemma \ref{L3.5} that \( (u_n) \) is bounded. Thus, we may assume (up to a subsequence) that \( u_n \rightharpoonup u^0 \).

We proceed by contradiction and assume that \( u^0 \equiv 0 \). 
Then it follows from Lemma \ref{L3.12} that
$$
\int_{\mathbb{R}^N} \left|\nabla \left(u_n  -  u^1(\cdot - y_n^1)\right)\right|^2 
+ \left|G^{-1}\left(u_n  - \ u^1(\cdot - y_n^1)\right)\right|^2\to 0.
$$
It follows from Lemma \ref{L3.13}  that $\|u_n-u^1(\cdot-y_n^1)\|\to0$, then 
$\|t_nv_n(\cdot+y_n^1)-u^1\|\to0$. Since $|\beta(t_nv_n(\cdot+y_n^1))|=|-y_n^1|\to\infty$, 
a contradiction.

It then follows from Lemmas \ref{L3.12} and \ref{L3.13} that \( u_n \to u^0 \) in \( E \). 
Arguing as in Proposition \ref{Pro1}, we finally obtain $u^0\in\mathcal{N}$.
\end{proof}

It remains to analyze the case $b > c_0 = c_\infty$. The proof of Lemma \ref{L3.11} ensures 
the existence of $S > 0$ and $R > 1$ such that for each $|y| \geq S$, one can find 
$T^y \in (0, R)$ satisfying $T^y u^{y,1}_\infty \in \mathcal{N}$. The map $y \mapsto T^y$ 
is continuous and satisfies $T^y \to 1$ as $|y| \to \infty$. However, we cannot guarantee  
that $u^{y,1}\in\mathcal{T}$ for $|y|< S$, and therefore a modification of our approach is necessary.

\begin{lemma}\label{L3.14} 
There exist $\theta_y\geq1$
and $T^y>0$, which are bounded and continuous with respect to $y$, 
such that $T^yu_\infty^{y,\theta_y}\in\mathcal{N}$ for every $y\in{\mathbb{R}^N}$. 
Moreover,  $\theta_y=1$ for all $|y|\geq S'+1$ with some $S'>0$.
\end{lemma}

\begin{proof}
We define $a:=\inf_{x\in{\mathbb{R}^N}}[s(x)-\frac{V(x)}{g^2(\infty)}]$, 
$a^*:=s_\infty-\frac{V_\infty}{g^2(\infty)}$ and 
$$
\mathcal{K}(y,\theta):=\frac{1}{\theta^2}\int_{\mathbb{R}^N}|\nabla u_\infty|^2-\int_{\mathbb{R}^N}[s(\theta x+y)-\frac{V(\theta x+y)}{g^2(\infty)}]u_\infty^2.
$$
It follows from (A5) that
$$
\mathcal{K}(y,\theta)\leq\frac{1}{\theta^2}\int_{\mathbb{R}^N}|\nabla u_\infty|^2-a\int_{\mathbb{R}^N}u_\infty^2,
$$
with equality if and only if $s(x)-\frac{V(x)}{g^2(\infty)}\equiv a$, then $a=a^*$. 
Given any  $c > 1$, let 
$\theta_0:=\max\{\sqrt{\frac{c}{a}}\cdot\frac{|\nabla u_\infty|_2}{|u_\infty|_2},1\}$, then 
$\mathcal{K}(y,\theta_0)\leq(\frac{1}{c}-1)a\int_{\mathbb{R}^N}u_\infty^2<0$.

Using (A3), (A5), (A7) and Lemma \ref{L2.1}-(4), we have 
\begin{align*}
\frac{\langle J'(ru_{\infty}^{y,\theta_0}),ru_{\infty}^{y,\theta_0}\rangle}{\theta^N_0r^{2}} 
&=\frac{1}{\theta^2_0}\int_{\mathbb{R}^{N}} \vert \nabla u_{\infty} \vert^{2} + \int_{\mathbb{R}^{N}} \frac{V(\theta_0 x+y) G^{-1}(ru_{\infty})}{g(G^{-1}(ru_{\infty}))ru_{\infty}} u_{\infty}^{2} \\
 &\quad  - \int_{\mathbb{R}^{N}}  \frac{f(\theta_0 x+y,G^{-1}(ru_{\infty}))}{g(G^{-1}(ru_{\infty}))ru_{\infty}} u_{\infty}^{2} \\
&=\mathcal{K}(y,\theta_0)+\int_{\mathbb{R}^{N}} V(\theta_0 x+y)[\frac{ G^{-1}(ru_{\infty})}{g(G^{-1}(ru_{\infty}))ru_{\infty}}-\frac{1}{g^2(\infty)}]u_{\infty}^{2}\\
&\quad  +\int_{\mathbb{R}^{N}} [s(\theta_0 x+y)- \frac{f (\theta_0 x+y,G^{-1}(ru_{\infty}))}{g(G^{-1}(ru_{\infty}))ru_{\infty}}]u_{\infty}^{2}\\
&\leq(\frac{1}{c}-1)a\int_{\mathbb{R}^N}u_\infty^2+V_\infty\int_{\mathbb{R}^{N}} [\frac{ G^{-1}(ru_{\infty})}{g(G^{-1}(ru_{\infty}))ru_{\infty}}-\frac{1}{g^2(\infty)}]u_{\infty}^{2} \\
&\quad  +\int_{\mathbb{R}^{N}} [s_\infty- \frac{f_\infty(G^{-1}(ru_{\infty}))}{g(G^{-1}(ru_{\infty}))ru_{\infty}}]u_{\infty}^{2}.
\end{align*}
Let 
$$
F_1(r)=\int_{\mathbb{R}^{N}} [\frac{ V_\infty G^{-1}(ru_{\infty})}{g(G^{-1}(ru_{\infty}))ru_{\infty}}-\frac{V_\infty}{g^2(\infty)}+s_\infty- \frac{f_\infty(G^{-1}(ru_{\infty})}{g(G^{-1}(ru_{\infty}))ru_{\infty}}]u_{\infty}^{2}.
$$ 
Hence, by a method analogous to Lemma \ref{L3.2}-(1), the function \(F_1(r)\) is strictly 
decreasing for \(r > 0\) and $\lim_{r\to\infty}F_1(r)=0$. 
Note that $\lim_{r\to0}F_1(r)\geq s_\infty\int_{\mathbb{R}^{N}}u_{\infty}^{2}$, 
then there is one and only one  $R_1>0$  for which 
$F_1(R_1)=(1-\frac{1}{c})a\int_{\mathbb{R}^N}u_\infty^2$, i.e., 
$R_1=F_1^{-1}((1-\frac{1}{c})a\int_{\mathbb{R}^N}u_\infty^2)$. 
Thus, $\frac{\langle J'(ru_{\infty}^{y,\theta_0}),ru_{\infty}^{y,\theta_0}\rangle}{\theta^N_0r^{2}}<0$, 
for $r>R_1$, uniformly for $y\in{\mathbb{R}^N}$. Similarly, we have $\frac{\langle J'(ru_{\infty}^{y,\theta_0}),ru_{\infty}^{y,\theta_0}\rangle}{\theta^N_0r^{2}}>0$ for $r$ small enough and
$\frac{\langle J'(ru_{\infty}^{y,\theta_0}),ru_{\infty}^{y,\theta_0}\rangle}{\theta^N_0r^{2}}$ 
is strictly decreasing on $r\in(0,\infty)$. Hence, there exists a unique $0<T^y\leq R_1$ 
such that $\frac{\langle J'(T^yu_{\infty}^{y,\theta_0}),T^yu_{\infty}^{y,\theta_0}\rangle}{\theta^N_0(T^y)^{2}}=0$, i.e., 
$T^y u_{\infty}^{y,\theta_0}\in\mathcal{N}$, for any $y\in{\mathbb{R}^N}$.

The argument in Lemma \ref{L3.11} yields 
\begin{align*}
&\frac{\langle J_\infty'(ru_{\infty}^{y,\theta}),ru_{\infty}^{y,\theta}\rangle}{\theta^Nr^{2}} \\
&=\frac{1}{\theta^2}\int_{\mathbb{R}^{N}} \vert \nabla u_{\infty} \vert^{2} + \int_{\mathbb{R}^{N}} \frac{V_\infty G^{-1}(ru_{\infty})}{g(G^{-1}(ru_{\infty}))ru_{\infty}} u_{\infty}^{2} 
  - \int_{\mathbb{R}^{N}}  \frac{f_\infty (G^{-1}(ru_{\infty}))}{g(G^{-1}(ru_{\infty}))ru_{\infty}} u_{\infty}^{2}\leq\alpha<0,
\end{align*} 
for $r>R$ where $\alpha$ and $R$ are as in Lemma \ref{L3.11}, uniformly for $\theta\geq1$. 
Then 
\begin{align*}
\frac{\langle J'(ru_{\infty}^{y,\theta}),ru_{\infty}^{y,\theta}\rangle}{\theta^Nr^{2}} 
&=\frac{\langle J_\infty'(ru_{\infty}^{y,\theta}),ru_{\infty}^{y,\theta}\rangle}{\theta^Nr^{2}}
 + \int_{\mathbb{R}^{N}} \frac{[V(\theta x+y)-V_\infty] G^{-1}(ru_{\infty})}{g(G^{-1}(ru_{\infty}))ru_{\infty}} u_{\infty}^{2} \\
&\quad + \int_{\mathbb{R}^{N}}  \frac{[f_\infty( G^{-1}(ru_{\infty}))-f(\theta x+y,G^{-1}(ru_{\infty}))]}{g(G^{-1}(ru_{\infty}))ru_{\infty}} u_{\infty}^{2}\\
&\leq\frac{\langle J_\infty'(ru_{\infty}^{y,\theta}),ru_{\infty}^{y,\theta}\rangle}{\theta^Nr^{2}}+\int_{\mathbb{R}^{N}}[s_\infty-s(\theta x+y)]u_{\infty}^{2}.
\end{align*} 
Using the Lebesgue dominated convergence theorem, we obtain a constant  $S'>0$ such that 
$\frac{\langle J'(ru_{\infty}^{y,\theta}),ru_{\infty}^{y,\theta}\rangle}{\theta^Nr^{2}} \leq\frac{\alpha}{2}<0$ 
for  $|y|\geq S'$ and $r>R$,  uniformly for $1\leq\theta\leq\theta_0$. 
Similarly, we have $\frac{\langle J'(ru_{\infty}^{y,\theta}),ru_{\infty}^{y,\theta}\rangle}{\theta^Nr^{2}}$ 
is strictly decreasing on $r\in(0,\infty)$ and 
$\frac{\langle J'(ru_{\infty}^{y,\theta}),ru_{\infty}^{y,\theta}\rangle}{\theta^Nr^{2}}>0$ for $r$ 
small enough, uniformly for $1\leq\theta\leq\theta_0$. 

Set $\theta_y:=\theta_0$, if $|y|<S'$; $\theta_y:=|y|-S'+\theta_0(S'+1-|y|)$, if 
$S'\leq|y|\leq S'+1$;
$\theta_y:=1$, if $|y|>S'+1$. So, there exist a unique $0<T^y\leq R_2:=\max \{R_1,R\}$ 
such that $\frac{\langle J'(T^yu_{\infty}^{y,\theta_y}),T^yu_{\infty}^{y,\theta_y}\rangle}{\theta^N_y(T^y)^{2}}=0$, i.e., $T^y u_{\infty}^{y,\theta_y}\in\mathcal{N}$, for any $y\in{\mathbb{R}^N}$. (Because the choice of $S'$ does not affect the proof of the subsequent lemmas, we can 
arbitrarily choose $\alpha > 0$ sufficiently small, i.e., $R > 1$ sufficiently small in the following.)

Obviously, $\theta_y$ is continuous with respect to $y$. 
Assume $y_n\to y^0$ as $n\to\infty$. Since $(T^{y_n})$ is bounded, by extracting a suitable subsequence, we find $T^{y_n}\to T^0$. Then 
$$
\langle J'(T^0u_{\infty}^{y^0,\theta_{y^0}}),T^0u_{\infty}^{y^0,\theta_{y^0}}\rangle=\lim_{n\to\infty}\langle J'(T^{y_n}u_{\infty}^{y_n,\theta_{y_n}}),T^{y_n}u_{\infty}^{y_n,\theta_{y_n}}\rangle=0.
$$ 
By the uniqueness of $T^{y^0}$, we have $T^{y^0}=T^0$. So, we obtain that $T^y$  depends continuously on $y$ from Heine theorem.
\end{proof}

We set \(\Gamma[y] := T^y u_{\infty}^{y,\theta_y}\) for \(y \in \mathbb{R}^N\). 
This defines a continuous operator \(\Gamma: \mathbb{R}^N \to \mathcal{N}\). 
The properties of $\beta$ implies
\begin{equation}
\label{q8}\beta(\Gamma[y])=y, 
\end{equation} 
which is given by \cite[Lemma 3.13 ]{Lehrer}.

\begin{lemma}\label{L3.15} 
$J(\Gamma[y])\to c_\infty$, as $|y|\to\infty$.
\end{lemma}

\begin{proof}
Since the limiting functional $J_\infty$ is translation-invariant, we have 
$J(\Gamma[y]) \to J_\infty(u_\infty) = c_\infty$.
\end{proof}

First, define $c_*:=\inf\{c>c_\infty: c$ is a critical value of $J_\infty\}$. 
Then, set $\tilde{c}:=\min\{c_*,2c_\infty\}$. From (A12) we obtain 
$c_\infty<\tilde{c}\leq 2c_\infty$.

\begin{lemma}\label{L3.16} 
Let $J^c:=\{u\in E:J(u)\leq c\}$. There exists $\delta>0$  satisfying 
$c_\infty+\delta<\min\{b,\tilde{c}\}$ such that $\beta(u)\neq0$ holds for
all $u\in\mathcal{N}\cap J^{c_\infty+\delta}$. 
\end{lemma}

The definition of $b$ immediately yields the conclusion of the above lemma.

\begin{lemma}\label{L3.17} 
Set $\mu(a):=\frac{2\theta^{-N}_0\tilde{c}-2c_\infty}{R^2_2|u_\infty|^2_2}$. 
Under this definition, hypothesis {\rm (A11)} is well-defined.
\end{lemma}

\begin{proof}
From Lemma \ref{L3.14}, we obtain that  
$R_2=\max\{F_1^{-1}((1-\frac{1}{c})a|u_\infty|^2_2), R\}$ and 
$\theta_0:=\max\{\sqrt{\frac{c}{a}}\cdot\frac{|\nabla u_\infty|_2}{|u_\infty|_2},1\}$.
Then (A11) implies that 
$a> \frac{|\nabla u_\infty|^2_2}{|u_\infty|^2_2}\cdot(\frac{c_\infty}{\tilde{c}})^{\frac{2}{N}}$. 
Define $c(a):=\frac{a| u_\infty|^2_2}{|\nabla u_\infty|^2_2}
\cdot(\frac{\tilde{c}}{c_\infty})^{\frac{2}{N}}$, 
for given $1<c\leq c(a)$, set
$$
F^c(s):=a^*-s-\frac{2{[\min\{\sqrt{\frac{s}{c}}\cdot\frac{| u_\infty|_2}{|\nabla u_\infty|_2},1\}}]^N\cdot\tilde{c}-2c_\infty}{[\max\{F_1^{-1}((1-\frac{1}{c})s|u_\infty|^2_2),R\}]^2
|u_\infty|^2_2}.
$$
Clearly, $F^c(a^*)\leq0$, $\lim_{s\to0}F^c(s)=a^*>0$ and $F^c(s)$ is strictly decreasing 
on $s\in(0,a^*]$. Thus, there exist a unique $0<a^c\leq a^*$ such that $F^c(a^c)=0$. 
Since we can arbitrarily  choose $R>1$ sufficiently
small, then (A11) holds exactly when there exists $1<c\leq c(a)$ such that 
$a\geq a^c$, i.e., $ a=a^*$ or 
\begin{align*} 
2{[\min\{\sqrt{\frac{a}{c}}\cdot\frac{| u_\infty|_2}{|\nabla u_\infty|_2},1\}}]^N\cdot\tilde{c}-2c_\infty>(a^*-a)|u_\infty|^2_2,
\end{align*} 
and 
$$
F_1\bigg(\sqrt{\frac{2{[\min\{\sqrt{\frac{a}{c}}\cdot\frac{| u_\infty|_2}{|\nabla u_\infty|_2},1\}}]^N\cdot\tilde{c}-2c_\infty}{(a^*-a)|u_\infty|^2_2}}\bigg)\leq(1-\frac{1}{c})a|u_\infty|^2_2.
$$

Note that $a^c\to a^*$ as $c\to1$, then (A11) is equivalent to $a\geq \inf_{1<c\leq C(a)}\{a^c\}$. 
Condition (A11) is monotonic in the sense that if it holds for some $a_1$, then it holds for all $a\in[a_1, a^*]$. Hence, there exists 
$a^0\in(\frac{|\nabla u_\infty|^2_2}{|u_\infty|^2_2}\cdot(\frac{c_\infty}{\tilde{c}})^{\frac{2}{N}}, a^*)$ 
such that (A11) holds if and only if $a^0\leq a\leq a^*$ (in fact, $a^0$ is the solution of 
$a=\inf_{1<c\leq C(a)}\{a^c\}$). In hypothesis (A11), we  choose $c=c_0$ such that 
$a^{c_0}=\inf_{1<c\leq C(a)}\{a^c\}$.
\end{proof}

\begin{lemma}\label{L3.18}
 If {\rm (A11) } holds, then $J(\Gamma[y])<\tilde{c}$.
\end{lemma}

\begin{proof} 
Using (A7), (A11) and Lemma \ref{L2.1}-(3), (4), we find that
\begin{align*}
J(\Gamma[y])
&=J_\infty(\Gamma[y])+\frac{1}{2}\int_{\mathbb{R}^N}[V(x)-V_\infty]|G^{-1}(T^yu_\infty^{y,\theta_y})|^2\\
&\quad +\int_{\mathbb{R}^N}F_\infty(G^{-1}(T^yu_\infty^{y,\theta_y}))-F(x,G^{-1}(T^yu_\infty^{y,\theta_y}))\\
&\leq\theta_0^{N}c_\infty+\frac{1}{2}\int_{\mathbb{R}^N}[\frac{V(x)}{g^2(\infty)}-\frac{V_\infty}{g^2(\infty)}](T^yu_\infty^{y,\theta_y})^2
  +\frac{1}{2}\int_{\mathbb{R}^N}[s_\infty-s(x)](T^yu_\infty^{y,\theta_y})^2\\
&\leq\theta_0^{N}c_\infty+\frac{\theta^{-N}_0\tilde{c}-c_\infty}{R^2_2|u_\infty|^2_2}\int_{\mathbb{R}^N}(T^yu_\infty^{y,\theta_y})^2\leq\tilde{c}.
\end{align*}
We claim that $J(\Gamma[y])<\tilde{c}$. Assume, for a contradiction, that there is some 
$y_0\in{\mathbb{R}^N}$ such that equality holds in all three inequalities above. 
The equality in the first relation implies that $\theta_0=1$. From the third equality holding, 
we can conclude that $T^{y_0}=R_2$, i.e., $T^{y_0}=R_1\geq R$ which leads to 
$s(x)-\frac{V(x)}{g^2(\infty)}\equiv a$ by the proof of Lemma \ref{L3.14}. Thus,
\begin{align*}
J(\Gamma[{y_0}])
&=\theta_0^{N}c_\infty+\frac{1}{2}\int_{\mathbb{R}^N}[\frac{V(x)}{g^2(\infty)}-\frac{V_\infty}{g^2(\infty)}](T^{y_0}u_\infty^{{y_0},\theta_{y_0}})^2
 +\frac{1}{2}\int_{\mathbb{R}^N}[s_\infty-s(x)](T^yu_\infty^{{y_0},\theta_{y_0}})^2\\
&=c_\infty<\tilde{c},
\end{align*} 
a contradiction.
\end{proof}

\begin{proof}[Proof of Theorem \ref{thm1.1}]
 Lemma \ref{L3.11} implies \(c_0 \leq c_\infty\). In the case where \(c_0 < c_\infty\), Proposition \ref{Pro1} directly yields a ground state solution of \eqref{new}. 
 Also, if  $b = c_0 = c_\infty$, the same conclusion follows  readily from Proposition \ref{Pro2}. Therefore, it remains to consider that $b>c_0=c_\infty$.

Assume by contradiction that the functional $J$  admits no critical value in the interval 
$(c_\infty ,\tilde{c})$.  Applying the deformation lemma (cf. \cite[p.86]{Struwe}), 
we obtain a deformation \(\eta\) mapping 
\(\Upsilon^{\tilde{c}} \setminus K_{\tilde{c}}\) into \(\Upsilon^{c_\infty + \delta}\). 
Here, \(\Upsilon^{c} := \{ w \in \mathcal{N}_0 : \Upsilon(w) \leq c \}\) and \(K_c\) denotes 
the set of all critical points $v$ of \(\Upsilon\) with \(\Upsilon(v)=c\), while \(\delta\) 
is as in Lemma \ref{L3.16}.  By Lemmas \ref{L3.7} and \ref{L3.10}, the flow \(\eta\) does not r
each the boundary of \(\mathcal{N}_0\), and it fixes every point in 
\(\Upsilon^{c_\infty + \delta}\), i.e., \(\eta(u) = u\) for all \(u \in \Upsilon^{c_\infty + \delta}\).

Lemmas \ref{L3.15} and \ref{L3.16} imply the existence of \(\rho_1 \geq S' > 0\) with the 
property that for any \(\rho > \rho_1\),
$$
c_\infty<\max_{|y|=\rho}J(\Gamma[y])=\max_{|y|=\rho}\Upsilon(n^{-1}(\Gamma[y]))<c_\infty+\delta<b.
$$
Let $\zeta: B_\rho(0)\mapsto \partial B_\rho(0)$ be a continuous map defined by
$$
\zeta(y):=\rho\cdot\frac{\beta\circ n\circ \eta \circ n^{-1}(\Gamma[y])}{|\beta\circ n\circ \eta \circ n^{-1}(\Gamma[y])|}.
$$ 
From \eqref{q8}, we have $\zeta(y)=y$ for all $y\in\partial B_\rho(0)$, which contradicts 
 \cite[D.11]{Willem}. Furthermore, following the argument in 
 \cite[Lemma 2.4]{Clapp}, we obtain that $u^0$  is sign-definite.  
 Then, by elliptic regularity theory, we deduce that $u^0\in C^{1,\alpha}({\mathbb{R}^N})$. 
 Finally, the maximum principle guarantees that  $u^0$ is positive. 
 \end{proof}

\subsection*{Acknowledgments}
The authors thank the  anonymous referees for  their thorough review and constructive feedback. 


\begin{thebibliography}{10}

\bibitem{Ambrosetti} A. Ambrosetti, G. Cerami, D. Ruiz;
Solitons of linearly coupled systems of semilinear non-autonomous equations on ${\mathbb{R}^N}$, 
\emph{J. Funct. Anal. } \textbf{254} (2008), 2816–2845.

\bibitem{Berestycki2} H. Berestycki, T. Gallou\"et, O. Kavian;
Equations de Champs scalaires euclidiens non linéaires dans le plan,  
\emph{C.R. Acad. Sci. Paris Ser. I Math.} \textbf{297} (1983), 307–310.

\bibitem{Berestycki} H. Berestycki, P. L. Lions;
Nonlinear scalar field equations, I. Existence of a ground state, 
\emph{Arch. Rational Mech. Anal.} \textbf{82}  (1983), 313–345. 

\bibitem{Carriao} P. C. Carri\~ao, R. Lehrer, O. H. Miyagaki;
Existence of solutions to a class of asymptotically linear Schr\"odinger equations in ${\mathbb{R}^N} $ via the Pohozaev manifold, 
\emph{J. Math. Anal. Appl.} \textbf{428} (2015), 165–183. 

\bibitem{Cerami} G. Cerami, D. Passaseo;
The effect of concentrating potentials in some singularly perturbed problems, 
\emph{Calc. Var.} \textbf{17} (2003), 257–281. 

\bibitem{Chang} K. C. Chang;
\emph{Methods in Nonlinear Analysis}, Springer-Verlag, Berlin, 2005.  

\bibitem{Clapp} M. Clapp, L. A. Maia;
A positive bound state for an asymptotically linear or superlinear Schr\"odinger equation, 
\emph{J. Diff. Eq.}  \textbf{260} (2016), 3173–3192. DOI: 10.1016/j.jde.2015.09.059

\bibitem{colin} M. Colin, L. Jeanjean;
 Solutions for a quasilinear Schr\"odinger equation: a dual approach, 
 \emph{Nonl. Anal.} \textbf{56} (2004), 213--226.

\bibitem{deng2} Y. B. Deng, S. J. Peng, J. X. Wang;
Nodal soliton solutions for generalized quasilinear Schr\"odinger equations, 
\emph{J. Math. Phys.} \textbf{55} (2014), 051501.

\bibitem{deng} Y. B. Deng, S. J. Peng, S. S. Yan;
Positive soliton solutions for generalized quasilinear Schr\"odinger equations with critical growth, 
\emph{J. Diff. Eq.} \textbf{258} (2015), 115--147. DOI: 10.1016/j.jde.2014.09.006

\bibitem{fang3} X. D. Fang;
A positive solution for an asymptotically cubic quasilinear Schr\"odinger equation, 
\emph{Commun. Pure Appl. Anal.} \textbf{18} (2019), 51--64.

\bibitem{fang1} X. D. Fang, Z. Q. Han;
 Existence of a Ground State Solution for a Quasilinear Schr\"odinger equation, 
 \emph{Adv. Nonlinear Stud.} \textbf{14} (2014), 941–950.

\bibitem{fang2} X. D. Fang, Z. Q. Han;
Ground state solutions for generalized quasilinear Schr\"odinger equations, 
\emph{Asymptot. Anal.} \textbf{140} (2024), 109--122.

\bibitem{Jeanjean} L. Jeanjean, T. J. Luo, Z. Q. Wang;
Multiple normalized solutions for quasi-linear Schr\"odinger equations, 
\emph{J. Diff. Eq.} \textbf{259} (2015), 3894–3928

\bibitem{Lehrer} R. Lehrer, L. A. Maia;
Positive solutions of asymptotically linear equations via Pohozaev manifold, 
\emph{J. Funct. Anal.} \textbf{266}  (2014), 213--246.

\bibitem{Wang1} J. Q. Liu, Z. Q. Wang;
Soliton solutions for quasilinear Schr\"odinger equations, I, 
\emph{Proc. Amer. Math. Soc.} \textbf{131} (2003), 441--448.

\bibitem{Wang2} J. Q. Liu, Y. Q. Wang, Z. Q. Wang;
Soliton solutions for quasilinear Schr\"odinger equations, II, 
\emph{J. Diff. Eq.} \textbf{187} (2003), 473--493.

\bibitem{Poppenberg} M. Poppenberg, K. Schmitt, Z. Q. Wang;
On the existence of soliton solutions to quasilinear Schr\"odinger equations, 
\emph{Calc. Var.} \textbf{14} (2002), 329--344.

\bibitem{shen} Y. T. Shen, Y. J. Wang;
Soliton solutions for generalized quasilinear Schr\"odinger equations, 
\emph{Nonl. Anal.} \textbf{80} (2013), 194--201. DOI: 10.1016/j.na.2012.10.005

\bibitem{Silva} E. A. B. Silva, G. F. Vieira;
Quasilinear asymptotically periodic Schr\"odinger equations with critical growth, 
\emph{Calc. Var. Partial Differ. Equ.} \textbf{39} (2010), 1–33.

\bibitem{Struwe} M. Struwe;
\emph{Variational Methods}, second ed., Springer-Verlag, Berlin, 1996. DOI: 10.1007/978-3-662-03212-1

\bibitem{Stuart} C. A. Stuart,;
An introduction to elliptic equation on ${\mathbb{R}^N}$, in \emph{Nonlinear Functional Analysis and Applications to Differential Equations} (A. Ambrosetti, K.C. Chang and I. Ekeland eds.), 
World Scientific, Singapore, 1998.

\bibitem{Szulkin1} A. Szulkin, T. Weth;
Ground state solutions for some indefinite variational problems, 
\emph{J. Funct. Anal.} \textbf{257}  (2009), 3802--3822. DOI: 10.1016/j.jfa.2009.09.013. 

\bibitem{Szulkin2} A. Szulkin,T. Weth;
The method of Nehari manifold, in \emph{Handbook of Nonconvex Analysis and Applications}, 
Int. Press, Boston, (2010), 597–632. 

\bibitem{Willem} M. Willem;
\emph{Minimax Theorems}, in \emph{Progress in Nonlinear Differential Equations and Their Applications, 24}, 
Birkhäuser Boston, Inc., Boston, (1996), x--162. DOI: 10.1007/978-1-4612-4146-1

\bibitem{Xue} Y. Xue, Z. Q. Han;
 Existence and multiplicity of solutions for Schr\"odinger equations with sublinear nonlinearities, 
 \emph{IMS Math.} \textbf{6} (2021), 5479--5492.

\bibitem{Yang} M. B. Yang,Y. H. Ding;
Existence of semiclassical states for a quasilinear Schr\"odinger equation with critical 
exponent in ${\mathbb{R}^N}$, 
\emph{Ann. Mat. Pura Appl.} \textbf{192} (2013), 783--804. DOI: 10.1007/s10231-011-0246-6

\end{thebibliography}

\end{document}
